\documentclass{article}
\usepackage[margin=1in]{geometry}
\usepackage{amsmath,amssymb}
\usepackage[most]{tcolorbox}
\usepackage[hidelinks]{hyperref}

\newcommand{\BookName}{Pattern Recognition and Machine Learning}
\newcommand{\BookAuthor}{Christopher M. Bishop}
\newcommand{\ChapterName}{Introduction}
\newcommand{\ExerciseID}{1.14}

\title{Chapter: \ChapterName}
\author{Ahonakpon Guy Gbaguidi}
\date{}

\begin{document}

\maketitle

\begin{center}
  \large\textbf{Exercise \ExerciseID}

  \vspace{0.5em}

  \normalsize\BookName\ - \BookAuthor
\end{center}

\begin{tcolorbox}[breakable,colback=gray!5,colframe=gray!60,title=Solution]
For any real square matrix $\mathbf{W}$, define
\[
\mathbf{W}^{S}=\frac{\mathbf{W}+\mathbf{W}^{\mathsf{T}}}{2},
\qquad
\mathbf{W}^{A}=\frac{\mathbf{W}-\mathbf{W}^{\mathsf{T}}}{2}.
\]
In terms of their elements,
\[
w_{ij}^{S}=\frac{w_{ij}+w_{ji}}{2},
\qquad
w_{ij}^{A}=\frac{w_{ij}-w_{ji}}{2}.
\]
These definitions give
\begin{align*}
w_{ji}^{S}
&=\frac{w_{ji}+w_{ij}}{2}=w_{ij}^{S}, \\
w_{ji}^{A}
&=\frac{w_{ji}-w_{ij}}{2}=-w_{ij}^{A}, \\
w_{ij}^{S}+w_{ij}^{A}
&=\frac{w_{ij}+w_{ji}}{2}+\frac{w_{ij}-w_{ji}}{2}=w_{ij}.
\end{align*}
Thus, $\mathbf{W}^{S}$ is symmetric, $\mathbf{W}^{A}$ is antisymmetric,
and $\mathbf{W}=\mathbf{W}^{S}+\mathbf{W}^{A}$. In particular,
$w_{ii}^{A}=0$ for every diagonal element.

Consider the contribution of the antisymmetric part to the quadratic form:
\begin{align*}
\sum_{i=1}^{D}\sum_{j=1}^{D} w_{ij}^{A}x_i x_j
&= \sum_{i=1}^{D}\sum_{j=1}^{D} w_{ji}^{A}x_j x_i
  \quad \text{(rename the dummy indices $i$ and $j$)} \\
&= -\sum_{i=1}^{D}\sum_{j=1}^{D} w_{ij}^{A}x_i x_j
  \quad \text{(antisymmetry and $x_i x_j=x_j x_i$)}.
\end{align*}

Let the original double sum be $A$. The equation above gives $A=-A$, so
$2A=0$ and therefore
\[
\sum_{i=1}^{D}\sum_{j=1}^{D} w_{ij}^{A}x_i x_j=0.
\]
Thus, only the symmetric part contributes:
\begin{align*}
\sum_{i=1}^{D}\sum_{j=1}^{D} w_{ij}x_i x_j
&= \sum_{i=1}^{D}\sum_{j=1}^{D}
    \left(w_{ij}^{S}+w_{ij}^{A}\right)x_i x_j
  \quad \text{(matrix decomposition)} \\
&= \sum_{i=1}^{D}\sum_{j=1}^{D} w_{ij}^{S}x_i x_j
  \quad \text{(the antisymmetric contribution is zero)}.
\end{align*}

It remains to count the independent parameters in the symmetric matrix
$\mathbf{W}^{S}$. There are $D$ independent diagonal elements. For the
off-diagonal elements, symmetry gives $w_{ij}^{S}=w_{ji}^{S}$, so only one
element from each pair is independent. There are
$D(D-1)/2$ such pairs. Hence, the total number of independent parameters is
\[
D+\frac{D(D-1)}{2}
=\frac{D(D+1)}{2}.
\]
Therefore, the quadratic form can be represented using only
$D(D+1)/2$ independent coefficients. Explicitly,
\[
\sum_{i=1}^{D}\sum_{j=1}^{D} w_{ij}x_i x_j
=\sum_{i=1}^{D} w_{ii}^{S}x_i^2
+2\sum_{1\leq i<j\leq D} w_{ij}^{S}x_i x_j.
\]
The factor of $2$ accounts for the two equal off-diagonal contributions,
so the coefficients can be chosen symmetric without loss of generality.
\end{tcolorbox}

\end{document}
